\chapter{Working Capital Is the Price of Growth}
\chaptermark{Working Capital Is the Price of Growth}

Louis, the founder of a labor-force management company, describes the first
large contract as nine million dollars of annualized revenue. The customer
would pay invoices after thirty days. His company, however, had to fund payroll
and other costs as the work was performed. By his rough calculation, five weeks
of outflow required about one million dollars before the first payment could be
trusted to arrive.

The contract was evidence of demand and a threat to solvency at the same time.
Refusing it would surrender a route to scale. Accepting it without financing
could leave workers unpaid even if the customer eventually paid every invoice.
The contradiction is only apparent. A sale records an exchange; cash records a
date.

Working capital pays for the interval between those dates. It carries the
inventory bought before sale, the labor performed before collection, the tax
reserved before payment, the new employee trained before productivity, and the
capacity installed before utilization. Growth usually enlarges that interval
before it rewards the owner. That is why a company can be busy, profitable on
paper, and unable to meet payroll.

\section{Put every number in its own account}

The language becomes dangerous when one impressive number is allowed to stand
for all the others. Keep at least five records separate:

\begin{description}[leftmargin=2.5em,style=nextline]
\item[Revenue] is the price assigned to delivered or contractually earned work,
  according to the applicable accounting rules. It is not cash available to
  spend.
\item[Gross profit] is revenue less the direct cost required to produce that
  revenue. It must still support overhead, financing, tax, repair, and loss.
\item[Operating profit] includes the period's operating expenses. It may contain
  revenue not yet collected and omit future cash commitments.
\item[Cash] is money presently available to the relevant person or entity. Cash
  in a company is not automatically the owner's personal spending money.
\item[Balance-sheet capacity] is the stock of cash, receivables, inventory,
  obligations, debt, and equity that can absorb timing and loss. Its quality
  matters as much as its size.
\end{description}

A Houston operator recalls a year in which he reports earning about three
million dollars, owing roughly nine hundred fifty thousand dollars in tax, and
having only about ten thousand dollars in his personal account because money
remained inside the company. The entity and tax details are not clear enough to
generalize from those figures. The durable point is that company profit,
retained cash, tax liability, and personal liquidity can move in different
directions.

This separation also corrects a common slogan about return on every dollar.
Some spending has measurable incremental revenue; some protects safety,
quality, law, maintenance, privacy, or human dignity. The second category is
not waste because attribution is difficult. A sound budget names the obligation
and the evidence appropriate to it rather than inventing a sales return.

\section{Draw the cash calendar before accepting the sale}

For one order, contract, or growth step, place actual expected cash movements on
a weekly calendar. Let \(I_t\) be cash received in week \(t\) and \(O_t\) cash
paid. The cumulative funding deficit through week \(t\) is

\begin{equation}
F_t=\sum_{\tau=0}^{t}\bigl(O_\tau-I_\tau\bigr).
\end{equation}

If \(B\) is a reserve for delay, defects, refunds, rework, and forecast error,
the minimum planned growth carry is

\begin{equation}
W_{\mathrm{peak}}=\max_t F_t+B.
\end{equation}

Use cash dates, not invoice dates. Include deposits, material orders, freight,
payroll, payroll taxes, commissions, subcontractors, insurance, implementation,
quality control, support, returns, debt service, and tax. Then move collection
later, not earlier, in the stress case. A stated thirty-day term can become
forty-five or sixty days after approval, dispute, customer processing, or a
weekend.

For inventory businesses, the familiar cash-conversion approximation is

\begin{equation}
\begin{aligned}
\text{cash-cycle days}={}&\text{inventory days}\\
&+\text{receivable days}\\
&-\text{payable days}.
\end{aligned}
\end{equation}

It is a diagnostic, not a forecast. Averages conceal the largest order, a late
customer, seasonal inventory, deposits that must be refunded, and wages that
cannot wait. The weekly cash calendar exposes the actual peak.

An apparel founder in Las Vegas explains the same problem through a large
retailer order: an order for one hundred racks can be a gift and a curse when
the supplier can finance only fifty. A purchase order is demand evidence, not
cash, margin, collection, or production capacity. Before celebrating it, verify
unit cost, defect allowance, shipping, chargebacks, cancellation, acceptance,
payment date, supplier terms, and the cash needed if the buyer pays late.

Another interview describes an e-cigarette company that gave retailers a
reported forty-percent gross margin, refused consignment, and collected up
front. Those terms shifted inventory risk toward the retailer and improved the
supplier's cash cycle. They do not make the enterprise a model. The speaker's
health claims are unsupported, the product is addictive and harmful, and
commercially effective distribution can scale harm as readily as benefit.
Cash architecture must be judged together with product truth and public cost.

\section{Choose explicitly who carries the gap}

The timing gap never disappears. Someone carries it: the founder, customer,
supplier, worker, lender, factor, or investor. The legitimate question is which
arrangement fits the evidence and allocates risk without deception or coercion.

Customer deposits and milestone billing can align payment with committed work,
provided refund rights, deliverables, and custody are clear. Supplier terms can
match material payment to collection, but not by forcing a fragile supplier to
finance a stronger buyer indefinitely. Smaller batches, staged sites, rented
capacity, or a slower launch can reduce the peak rather than finance it. Retained
earnings preserve control but cost time and may be unavailable to founders who
did not begin with margin.

A revolving bank line may provide lower-cost short-term liquidity when the
business has credible records and repayment capacity. It also brings interest,
fees, covenants, security interests, renewal risk, and default consequences.
Equity has no scheduled repayment, but it transfers economics, information,
governance, and sometimes control. Neither is free simply because its price is
not visible on an invoice.

Louis used receivables factoring. He reports advancing about ninety percent of
eligible invoices at an all-in cost he describes as roughly twenty-one percent,
then retaining earnings until dependence on that facility fell. A factor can
turn a creditworthy invoice into earlier cash, but its full burden depends on
time, fees, reserve, recourse, disputes, dilution, customer concentration, and
who absorbs nonpayment. Compare total cash received and total cash surrendered
over the actual interval. Do not compare an unexplained fee with an annual bank
rate.

Expensive bridge finance can be rational when contribution after financing is
positive, collection is highly credible, the contract creates durable
capability, and the downside remains survivable. It can also preserve a contract
that should have been declined. The availability of money does not repair a
negative margin, weak customer, unsafe workload, or obligation the company
cannot deliver.

\section{Become bankable before the emergency}

A healthcare entrepreneur names a capable accountant as an early scaling hire
because trusted books support credit and an eventual sale. A Salt Lake City
property operator supplies the necessary counterweight: the owner cannot
delegate understanding. The accountant should produce reliable records and
challenge errors; the owner must still be able to explain revenue recognition,
gross margin, expenses, cash, receivables, payables, debt, tax reserves, and the
assumptions behind the forecast.

Bankability grows from evidence that another party can examine:

\begin{itemize}[leftmargin=1.5em,itemsep=0.15em]
\item reconciled accounts and timely statements;
\item contracts, invoices, aging, and collection history;
\item unit economics tied to actual delivery;
\item tax, payroll, insurance, and legal obligations kept current;
\item a cash forecast whose prior errors are visible;
\item owner capital, retained earnings, collateral, or other loss absorption;
\item a specific borrowing purpose and credible repayment source; and
\item early notice of adverse changes and covenant pressure.
\end{itemize}

A private-equity buyer in Palm Beach describes how an operating track record,
assets, lender trust, and business cash flow made acquisition borrowing
possible. These are not permission to pledge everything. Assets may be
overvalued and cash flow can fall. The debt must survive the stress case without
quietly transferring household ruin, unpaid labor, or customer loss to others.

During the collapse in travel after September 11, a hospitality founder recalls
his hotel cash registers going to zero. He contacted his banks repeatedly, disclosed
the good, bad, and ugly, and borrowed more to survive. Additional debt during a
zero-revenue shock can deepen insolvency; candor does not guarantee rescue. It
does, however, preserve the information a lender needs to consider a repair.
This directly contradicts another interviewee's advice to conceal parts of a
financial statement. In a lawful credit relationship, omission and
misrepresentation destroy the very trust on which flexibility depends.

\section{Stage growth so the company can learn}

An industrial-service operator in Palm Beach describes beginning with one
productive vacuum truck and expanding offices in stages from five to ten to
fifteen rather than attempting nationwide scale at once. The pace let retained
profit, operating knowledge, customers, and management capacity grow together.
It preserved control, though slower self-financing is not available or optimal
in every market.

Glenn Boyd reports that his fast-growing software company considered venture
capital, did not receive it, became more frugal, and later retained a very large
founder stake at its public offering. That is one outcome, not proof that
rejected financing causes success. His stronger warning is that expenses
multiply with growth and one adverse movement can threaten the whole company.

Margin and cash velocity must therefore be read together. If a business earns
small contribution per cycle but repeats a short, reliable cash cycle many
times, annual contribution may be substantial. A high stated margin tied up in
slow inventory or doubtful receivables may finance very little. Faster turnover
is valuable only if quality, returns, service cost, and collection remain
truthful. Speed cannot redeem a loss on every turn.

The same caution applies to emergency selling. A Scottsdale operator reports a
moment with a \$550,000 card bill, more than \$200,000 of payroll, and about
\$300,000 in cash, followed by roughly \$800,000 of new product sold in a week
and a half. The response reportedly closed his immediate gap. It could also
have enlarged it if cash had not been collected quickly or if delivery,
refunds, commissions, and support exceeded the contribution. Sales create
liquidity only to the extent that collected cash arrives before the obligations
created by the sale.

\section{Complete the growth-carry worksheet}

Choose one proposed contract, location, hire, product run, or acquisition. Do
not approve it from revenue alone.

\begin{enumerate}[itemsep=0.18em]
\item Record the demand evidence, cancellation rights, concentration, and
  realistic collection date.
\item Calculate unit revenue, direct cash cost, gross contribution, service and
  recovery cost, and the volume at which fixed capacity becomes necessary.
\item Draw weekly cash inflows and outflows through collection. Calculate
  \(W_{\mathrm{peak}}\) with a named reserve.
\item Name the production, quality, safety, support, and management capacity
  required at each stage. Remove work before adding an overloaded role.
\item Compare deposits, milestones, supplier terms, staging, retained earnings,
  bank credit, factoring, and equity by total cost, control, recourse, and
  failure consequence.
\item Run delay, lower-volume, lower-margin, defect, refund, and financing-loss
  cases. Include tax and owner household boundaries.
\item Set stop conditions before commitment: maximum uncovered deficit, latest
  collection date, minimum contribution, quality threshold, capacity limit,
  covenant headroom, and the event that pauses new sales.
\item Assign one person to update actual cash against forecast each week and
  explain every material variance without hiding it.
\end{enumerate}

Growth is earned when demand, contribution, cash timing, financing, and
capacity can travel together. If one cannot, the answer may be a higher price,
better term, smaller batch, later start, different capital, or a respectful no.
The spreadsheet can reveal the mismatch. It cannot carry the workload, absorb
the fear, challenge a bad promise, or decide whose safety matters. Those duties
belong to people, and the company now has to become worthy of theirs.
